\chapter{Conclusion}

In this dissertation I have explained how I, along with my co-authors, investigated various extensions of the
classical double copy.

To summarize our results, we found an extension of the Weyl double copy, which is formulated in terms of spinors,
to a large class of five dimensional spacetimes (the Kerr-NUT-(A)dS spacetimes) using an existing proposal on five
dimensional spinor representations. We additionally showed that the double copy structure for Kerr-NUT-(A)dS is
consistent with the double copy structure found using the Kerr-Schild double copy, as well as with existing
literature that used the Kerr-Schild double copy for various special cases within the Kerr-NUT-(A)dS class.

Next, in trying to understand the emergence of black hole horizons in the double copy paradigm, we found a way to
diagnose and locate apparent horizons in Kerr-Schild spacetimes. This involved noticing that for Kerr-Schild
spacetimes, the expansions of null geodesic congruences that are tangent to the Kerr-Schild vector field can be
rewritten using the single copy gauge field and the zeroth copy scalar. Thus, for stationary black holes that have
horizons generated by a Kerr-Schild vector field, in principle the location of their horizons can be found via
a calculation that uses only single and zeroth copy objects.

This approach is not without its drawbacks; for one, the calculation that we end up doing in the single and zeroth
copy does not respect the gauge freedom of the single copy. This does not by itself disprove our approach --- the
Kerr-Schild double copy is itself gauge dependent, and thus it is not a surprise that our proposal for finding
horizons is not gauge invariant. However the gauge dependence does mean that it is impossible to physically
interpret the quantities calculated if we were to restrict ourselves to a naive combination of the single and
zeroth copy theories, since all observable quantities should be gauge invariant.

Lastly, we turned to the question of the double copy structure of generic spacetimes. Our approach was to start
by restricting to a very special subset of spacetime (a null geodesic) and examining the double copy structure of
spacetime around the null geodesic to various orders of the scaling parameter. The lowest order corresponds to the
Penrose limit, where the spacetime metric (appropriately scaled) is that of a plane wave, whose double copy
structure is fairly well understood. We found that the procedure of taking the Penrose limit and of expanding about
it is consistent with the double copy. Moreover, we saw that for a generic spacetime, the existence of a local Weyl
double copy encounters an obstruction at second order beyond the Penrose limit.
